%%%%%%%%%%%%%%%%%%%%%%%%%%%%%%%%%%%%

\section{グラフを用いたモデルの条件確認}

%%%%%%%%%%%%%%%%%%%%%%%%%%%%%%%%%%%%

\subsection{診断プロット}

%%%%%%%%%%%%%%%%%%%%%%%%%%%%%%%%%%%%

\begin{frame}
\frametitle{モデルの条件（仮定）}

\[ \hat{y} = \beta_0 + \beta_1 x_1 + \beta_2 x_2 + \cdots + \beta_p x_p \]

$\:$ \\

モデルは以下の条件に依存する
\begin{enumerate}
\item 残差がほぼ正規分布に従う（大きなデータセットでは重要性が低い）
\item 残差の分散が一定である
\item 残差が独立である
\item 各変数が応答変数と線形の関係にある
\end{enumerate}

これらの条件の妥当性を確認するためにグラフによる方法をよく用いる。以下のスライドで詳しく解説する。

\end{frame}

%%%%%%%%%%%%%%%%%%%%%%%%%%%%%%%%%%

\begin{frame}
\frametitle{（1）ほぼ正規分布に従う残差}

残差の正規確率プロットおよび/またはヒストグラム：\\

\begin{center}
\includegraphics[width=0.6\textwidth]{9-3_model_cond/figures/beauty/normal_res}
\end{center}

\dq{この条件は満たされているように見えるか？}

\end{frame}

%%%%%%%%%%%%%%%%%%%%%%%%%%%%%%%%%%

\begin{frame}
\frametitle{（2）残差の分散が一定}

残差および/または残差の絶対値 vs. 当てはめ値（予測値）の散布図：\\

\begin{center}
\includegraphics[width=\textwidth]{9-3_model_cond/figures/beauty/homo_res}
\end{center}

\dq{この条件は満たされているように見えるか？}

\end{frame}

%%%%%%%%%%%%%%%%%%%%%%%%%%%%%%%%%%

\begin{frame}
\frametitle{一定分散の確認 --- まとめ}

\begin{itemize}

\item 単純線形回帰（1つの説明変数）では、\hl{残差 vs. $x$} のプロットを用いて一定分散の条件を確認した。

\item 重線形回帰（2つ以上の説明変数）では、\hl{残差 vs. 当てはめ値} のプロットを用いて一定分散の条件を確認する。

\end{itemize}

$\:$ \\

\dq{なぜ異なるプロットを使うのか？}

\soln{重線形回帰では説明変数が多いため、そのうちの1つに対する残差のプロットだけでは全体像を把握できないからである。
}

\end{frame}

%%%%%%%%%%%%%%%%%%%%%%%%%%%%%%%%%%

\begin{frame}[fragile]
\frametitle{（3）独立な残差}

データ収集順序に対する残差の散布図：\\

\begin{center}
\includegraphics[width=0.8\textwidth]{9-3_model_cond/figures/beauty/indep_res}
\end{center}

\dq{この条件は満たされているように見えるか？}

\end{frame}

%%%%%%%%%%%%%%%%%%%%%%%%%%%%%%%%%%

\begin{frame}
\frametitle{独立な残差の条件についてさらに}

\begin{itemize}

\item 残差が独立であることを確認することで、間接的に観測値の独立性を確認できる。

\item 観測値と残差が独立であれば、データ収集順序に対する残差の散布図に増加または減少トレンドは見られないはずである。

\item この条件は時系列データでしばしば満たされない。そのようなデータには適切な分析のために、より高度な時系列回帰技法が必要になる。

\end{itemize}

\end{frame}

%%%%%%%%%%%%%%%%%%%%%%%%%%%%%%%%%%

\begin{frame}[fragile]
\frametitle{（4）線形関係}

各（数値）説明変数に対する残差の散布図：

\begin{center}
\includegraphics[width=0.8\textwidth]{9-3_model_cond/figures/beauty/linear}
\end{center}

\dq{この条件は満たされているように見えるか？}

\Note{y 軸に予測変数の代わりに残差を使うのは、予測変数間の多重共線性などの他の違反を気にせずに線形性を確認できるからである。}

\end{frame}

%%%%%%%%%%%%%%%%%%%%%%%%%%%%%%%%%%%

\subsection{モデル適合を改善するためのオプション}

%%%%%%%%%%%%%%%%%%%%%%%%%%%%%%%%%%%%

\begin{frame}
\frametitle{モデルを改善するためのいくつかの方法}

\begin{itemize}

\item 変数の変換
\item モデルのギャップを埋める追加変数を探す
\item 分散の不一致や予測変数と応答変数の非線形関係などの課題に対処できる、より高度な手法を使用する

\end{itemize}

\end{frame}

%%%%%%%%%%%%%%%%%%%%%%%%%%%%%%%%%%%%

\begin{frame}
\frametitle{変換}

モデルの問題が説明変数と応答変数の間の非線形関係である場合、応答変数を変換することが有効なことがある。

\begin{itemize}

\item 対数変換（$\log y$）
\item 平方根変換（$\sqrt{y}$）
\item 逆数変換（$1/y$）
\item 切り捨て（最大値を制限する）

\end{itemize}

説明変数に変換を適用することも可能だが、そのような変換はモデル係数の解釈をさらに困難にする傾向がある。

\end{frame}

%%%%%%%%%%%%%%%%%%%%%%%%%%%%%%%%%%%%

\begin{frame}
\frametitle{モデルは間違っていても役に立つ}

\begin{quote}
すべてのモデルは間違っているが、いくつかは役に立つ。--- George Box
\end{quote}

\begin{itemize}

\item 完璧なモデルは存在しないが、不完全なモデルでもモデルの欠点を明確に報告する限り役に立つ可能性がある。

\item 条件が大きく違反している場合はモデルの結果を報告すべきではなく、より多くの統計的手法を学ぶか、専門家の助けを借りてでも新しいモデルを検討すべきである。

\end{itemize}

\end{frame}

%%%%%%%%%%%%%%%%%%%%%%%%%%%%%%%%%%%%
